\section{Problem setup}
\label{sec:problem}

\paragraph{The black box.} A trained classifier $f:\mathcal{X}\to[0,1]$ maps an
instance $x=(x_1,\dots,x_d)$ to a score, the estimated probability of the
favourable outcome (loan granted, credit line extended, low diabetes risk). A
threshold $\tau$ turns the score into a decision
$\hat y(x)=\mathbb{1}[f(x)\ge\tau]$. We treat $f$ as \emph{sacred}: the only
things any method --- ours included --- may ask of it are $f(x)$ and the induced
$\hat y(x)$. No gradients, no trees, no feature importances. This is both a
faithfulness requirement (real deployed models are often unavailable for
inspection) and a discipline that keeps the comparison honest, since the
baselines that can peek would otherwise be solving a different problem.

Fix an applicant $x^{0}$ who was rejected, $f(x^{0})<\tau$. Recourse is advice
that moves them to the favourable side.

\paragraph{Two things called ``a counterfactual.''} The literature and this paper
use the word for two different objects, and the gap between them is the whole
story.
\begin{itemize}
  \item A \emph{counterfactual point} $x'$ is an instance with
    $f(x')\ge\tau$. This is what distance-minimising methods return: a
    destination.
  \item An \emph{action} $a=\{(j,v_j):j\in S\}$, $S\subseteq\{1,\dots,d\}$,
    $|S|\le K$, is a set of interventions $\dopart{X_S=v_S}$ the applicant
    performs on features they control. This is a route.
\end{itemize}
A point is not a route. Setting $x'_j$ for the features a method chose to change
does not leave the other features where the method left them, because features
cause one another.

\paragraph{The world an action produces.} We equip each dataset with a structural
causal model $\mathcal{M}$ (\S\ref{sec:scm}): a directed acyclic graph over the
features and, for each feature, a structural equation
$X_j := g_j\big(\mathrm{pa}(j),U_j\big)$ from its parents and an exogenous term.
Given the applicant $x^{0}$ and an action $a$, the record that actually results is
obtained by the standard abduction--action--prediction steps
\citep{pearl2009causality}: infer the exogenous terms consistent with $x^{0}$,
apply $\dopart{X_S=v_S}$, and recompute every descendant of $S$ through its
structural equation. Write the result
\begin{equation}
  x^{a} \;=\; \mathcal{M}\big(x^{0};\,\dopart{X_S=v_S}\big).
  \label{eq:propagate}
\end{equation}
Features that are not descendants of $S$ keep their values; descendants move to
whatever the equations imply. It is $x^{a}$, not the named point, that the world
presents to the black box after the applicant acts.

\paragraph{The mismatch, stated plainly.} A point-based method reports success
when $f(x')\ge\tau$. But the applicant, acting on the changes in $x'$, arrives at
$x^{a}$, and the decision they actually receive is $\hat y(x^{a})$. When
$f(x')\ge\tau$ but $f(x^{a})<\tau$, the advice flipped the model on paper and
failed the person in fact. Section~\ref{sec:desiderata} turns this mismatch, and
the separate question of whether $a$ is something a person could even do, into
measurable audits; Section~\ref{sec:method} builds a method that scores $x^{a}$
directly and so cannot fall into the gap.
